\chapter{Formal Problem Definition}
\label{ch:formalization}

This chapter is the mathematical core of the report.
It does not yet describe a pipeline; it fixes the objects a pipeline must
manipulate.
The industrial-monitoring contract \code{classify\_signal} from
Chapter~\ref{ch:intro} is the sole running example: every definition is
instantiated on it before the next object is introduced.
Notation: $\mathbb{N}$, $\mathbb{Z}$, and $\mathbf{1}[\cdot]$ for the Iverson
bracket.

% ---------------------------------------------------------------------------
\section{Running Example Roadmap}
\label{sec:form:roadmap}
% ---------------------------------------------------------------------------

What must be true of \code{classify\_signal} before any synthesis loop is
built?
The answer is a short chain: a typed task, an ordered oracle, a finite witness
set that respects first-match priority, and a release predicate under budget~$K$.
Figure~\ref{fig:specir-walkthrough} previews that chain from FSF text to the
witness $(-3,-10)$.

\begin{figure}[t]
  \centering
  \tikzinput{diagrams/tikz/specir_walkthrough}
  \caption{\textbf{From FSF text to witness (classify\_signal).}
    The adapter lowers surface syntax into SpecIR; the witness generator
    encodes first-match priority as $\Phi_i$; Z3 returns a model such as
    $(-3,-10)$ that activates $S_1$.
    Takeaway: SpecIR is the hinge between notation-specific text and the
    shared prevention pipeline.}
  \label{fig:specir-walkthrough}
\end{figure}

% ---------------------------------------------------------------------------
\section{Formal Task Model}
\label{sec:form:task}
% ---------------------------------------------------------------------------

Release decisions need a precise unit of work: what must the pipeline know
about a sprint artefact before it can judge a candidate?
That unit is a \emph{formal specification task}---a single function whose
behaviour is described entirely by an ordered scenario specification.
HSP-Agile instantiates the model over SOFL/FSF~\citep{li2022agilesofl,liu2014sofl};
every later stage consumes one such task.

\begin{definition}[Formal Specification Task]
\label{def:task}
A \emph{formal specification task} is a triple
\[
  t \;=\; \bigl(\mathcal{S}_t,\;\mathcal{X}_t,\;\mathcal{Y}_t\bigr),
\]
where
\begin{itemize}
  \item $\mathcal{S}_t = \langle s_1, s_2, \ldots, s_n \rangle$ is a finite,
        ordered sequence of $n$ FSF scenarios (the \emph{specification});
  \item $\mathcal{X}_t \subseteq \mathbb{Z}^d$ is the \emph{input domain},
        a $d$-dimensional integer vector space; and
  \item $\mathcal{Y}_t \subseteq \mathbb{Z}^m$ is the \emph{output domain},
        an $m$-dimensional integer vector space.
\end{itemize}
\end{definition}

\paragraph{Example.}
For \code{classify\_signal}, $d=2$, $m=1$ (string-valued alert encoded as a
discrete label), $\mathcal{S}_t=\langle s_1,s_2,s_3\rangle$, and
$\mathcal{X}_t\subseteq\mathbb{Z}^2$ with coordinates
$(\mathit{level},\mathit{threshold})$.
Benchmark tasks in later chapters often use $d=5$ and $m=3$; the same
definitions apply.

% ---------------------------------------------------------------------------
\section{Specification Intermediate Representation (SpecIR)}
\label{sec:form:specir}
% ---------------------------------------------------------------------------

The prevention pipeline should not be rewritten for every surface notation.
What is needed is a shared hinge: a notation-agnostic form that still preserves
ordered guarded cases.
That hinge is the \textbf{Specification Intermediate Representation}
(\textbf{SpecIR}).
Adapters map SOFL/FSF, Mini-Z, or Mini-StateMachine text into SpecIR; the rest
of the pipeline depends only on SpecIR.

\begin{definition}[Specification Intermediate Representation]
\label{def:specir}
For task $t$, $\text{SpecIR}(t) = (\Sigma_t,\; \Gamma_t,\; C_t,\; M_t)$, where
$\Sigma_t$ is the typed I/O signature, $\Gamma_t$ the variable domain,
$C_t = \langle c_1,\ldots,c_n\rangle$ an ordered list of guarded cases
(first-match semantics), and $M_t$ notation-specific metadata.
The abstract syntax is written
$\text{SpecIR} \triangleq \langle\Sigma_t, \Gamma_t, C_t, M_t\rangle$ or,
equivalently, $\text{SpecIR} ::= \texttt{Signature}\;\texttt{Cases}\;[\texttt{Metadata}]$.
\end{definition}

\paragraph{Example.}
For \code{classify\_signal}, $\Sigma_t$ declares inputs
\texttt{level:int}, \texttt{threshold:int} and output \texttt{alert:string};
$C_t$ holds three guarded cases whose first guard atom is
$\{\texttt{var}{=}\texttt{level},\;\texttt{op}{=}\texttt{lt},\;
\texttt{threshold}{=}0\}$ (Figure~\ref{fig:specir-walkthrough}).

The abstract tuple is not yet checkable.
Adapters and validators need a concrete surface form that matches the JSON
schema used in the implementation.

\begin{definition}[SpecIR Surface Grammar (EBNF)]
\label{def:specir-ebnf}
The JSON serialisation aligns with \texttt{schemas/spec\_ir.schema.json} and
\texttt{src/ir/spec\_ir.py}:
\begin{lstlisting}[language={},basicstyle=\ttfamily\footnotesize]
SpecIR      ::= task_id notation name Signature Cases surface_prompt [Metadata]
Signature   ::= Inputs Outputs
Inputs      ::= Param*          Outputs ::= Param*
Param       ::= name type       (* nat | int | bool | string *)
Cases       ::= GuardedCase+
GuardedCase ::= index guard postcondition [guard_text] [post_text]
Guard       ::= GuardAtom* | { var="" op="others" }
GuardAtom   ::= var op [threshold]   (* op in eq,ne,lt,le,gt,ge,others *)
\end{lstlisting}
\end{definition}

Parsing alone is not enough: a document can parse and still omit variables,
skip indices, or leave residual inputs uncovered.
Well-formedness rules those cases out.

\begin{definition}[SpecIR Well-formedness]
\label{def:specir-wf}
A SpecIR document $d$ is \emph{well-formed} iff:
(i)~$d.\texttt{notation}$ is one of \texttt{sofl}, \texttt{mini\_z},
\texttt{mini\_statemachine}, or \texttt{real\_derived};
(ii)~case indices are consecutive $1,\ldots,n$ with pairwise distinct guards;
(iii)~every \texttt{guard} atom's \texttt{var} is declared in \texttt{inputs},
    except the \texttt{others} sentinel;
(iv)~every \texttt{postcondition} key is declared in \texttt{outputs}; and
(v)~the final case is either an \texttt{others} sentinel or its guard is
    logically equivalent to $\neg\bigvee_{i<n} g_i$.
Validation is enforced by \texttt{src/ir/schema\_validate.py}
(Appendix~\ref{sec:repro:schemas}).
\end{definition}

\paragraph{Why adapters and lowering exist.}
Surface notations still differ.
A thin map turns notation-specific text into well-formed SpecIR; after that
step the shared pipeline never sees the original syntax.
HSP-Agile optionally projects SpecIR back onto an FSF-shaped view for
screening that still speaks in scenarios.

\begin{definition}[Specification Adapter]
\label{def:spec-adapter}
A \emph{Specification Adapter} is $\text{Adapt}_L : \text{Text}_L \longrightarrow \text{SpecIR}$.
\end{definition}

On the SOFL path, screening still speaks in scenarios, so SpecIR is projected
onto an FSF-shaped \texttt{TaskSpec}; the core stages continue to depend only
on SpecIR.

\begin{definition}[FSF Lowering (HSP-Agile Instantiation)]
\label{def:fsf-lowering}
$\text{Lower}_{\text{FSF}} : \text{SpecIR} \longrightarrow \text{TaskSpec}_{\text{FSF}}$
projects SpecIR onto the FSF view; the pipeline core depends only on SpecIR.
\end{definition}

Each ordered SpecIR case is one FSF scenario: a guard that selects a region,
and an output mapping that says what must be returned there.

\begin{definition}[FSF Scenario]
\label{def:scenario}
A \emph{Functional Scenario Framework (FSF) scenario} is a pair
$s_i = (g_i,\;\phi_i)$, where $g_i : \mathcal{X}_t \to \{0, 1\}$ is a Boolean
guard and $\phi_i : \mathcal{X}_t \to \mathcal{Y}_t$ is a total output mapping.
\end{definition}

\paragraph{Example.}
$s_1=(\mathit{level}<0,\;\phi_1)$ with $\phi_1\equiv\texttt{"Error"}$;
$s_2=(\mathit{level}\ge\mathit{threshold},\;\phi_2)$ with
$\phi_2\equiv\texttt{"Critical"}$; $s_3$ is the \texttt{others} case with
$\phi_3\equiv\texttt{"Normal"}$.
Guards are linear arithmetic predicates over $\mathcal{X}_t$; each scenario is
a self-contained input--output contract.

Priority is not an implementation detail.
It is part of the specification, and the oracle is the function that enforces
it~\citep{dijkstra1975guarded}.

\begin{definition}[Scenario Ordering and Specification Oracle]
\label{def:oracle}
Scenario $s_i$ has \emph{higher priority} than $s_j$ whenever $i < j$.
For any input $x \in \mathcal{X}_t$, the \emph{specification oracle}
$g_t : \mathcal{X}_t \to \mathcal{Y}_t$ is
\[
  g_t(x) \;=\; \phi_i(x),
  \qquad \text{where } i = \min\bigl\{j : g_j(x) = 1\bigr\}.
\]
The final scenario $s_n$ is the \emph{``others'' case} with
$g_n(x) = \neg\bigvee_{j=1}^{n-1} g_j(x)$, so every input is handled by
exactly one scenario.
\end{definition}

\paragraph{Example.}
On $(-3,-10)$, both $g_1$ and $g_2$ hold, so
$g_t(-3,-10)=\phi_1(-3,-10)=\texttt{"Error"}$.
An implementation that checks $g_2$ first returns \texttt{"Critical"} and
fails the oracle---correct on most disjoint inputs, wrong wherever both
guards hold.

% ---------------------------------------------------------------------------
\section{FSF Scenario Semantics}
\label{sec:form:semantics}
% ---------------------------------------------------------------------------

The oracle in Definition~\ref{def:oracle} is a mathematical object.
The pipeline must make it operational: it needs concrete inputs on which to
execute candidates and compare outputs.
Those inputs are the concrete cases below, generated by SMT solving rather than
by random sampling.

\begin{definition}[Concrete Case]
\label{def:concrete-case}
A \emph{concrete case} for task $t$ is a pair
\[
  (x,\;y) \;=\; \bigl(x,\;g_t(x)\bigr),
  \qquad x \in \mathcal{X}_t.
\]
The set of all concrete cases is called the \emph{case set} $\mathcal{D}(t)$.
\end{definition}

In practice $\mathcal{X}_t$ is too large for exhaustive enumeration, so
$\mathcal{D}(t)$ is constructed by targeted SMT sampling \citep{demoura2008z3, barrett2018smt}:
one representative input per scenario.

\begin{definition}[SMT-Generated Test Suite]
\label{def:test-suite}
The case set $\mathcal{D}(t)$ is generated by the Z3 SMT solver as
follows.  For each scenario $s_i = (g_i, \phi_i)$, solve the satisfiability
query
\[
  x_i \;\models\; g_i(x) \;\land\; \bigwedge_{j < i} \neg\,g_j(x).
\]
If satisfiable, add $\bigl(x_i,\;g_t(x_i)\bigr)$ to $\mathcal{D}(t)$.
\end{definition}

\paragraph{Example.}
For $s_1$ the query is simply $\mathit{level}<0$; a model is $(-3,-10)$ with
oracle output \texttt{"Error"}.
For $s_2$ the query is
$\mathit{level}\ge\mathit{threshold}\land\neg(\mathit{level}<0)$; a model
might be $(5,3)$ with oracle output \texttt{"Critical"}.

The conjunct $\bigwedge_{j < i} \neg\,g_j(x)$ enforces strict scenario
membership: $x_i$ activates $s_i$ \emph{and no higher-priority scenario}.
This guarantees that the oracle maps $x_i$ through $\phi_i$ and not through
any $\phi_j$ with $j < i$, making each case a direct witness to the
prescribed output of exactly one scenario.

The following proposition establishes that the procedure achieves full
scenario coverage.

\begin{lemma}[Scenario Coverage]
\label{lem:coverage}
If the SMT solver finds satisfying inputs for all $n$ scenarios, then
$\mathcal{D}(t)$ contains at least one representative case per scenario,
covering the entire precedence ordering $s_1 \succ s_2 \succ \cdots \succ s_n$.
\end{lemma}

\begin{proof}
By construction, each $x_i$ satisfies $g_i$ but not $g_1, \ldots, g_{i-1}$.
Hence $i = \min\{j : g_j(x_i) = 1\}$, so the oracle evaluates
$g_t(x_i) = \phi_i(x_i)$.  Each distinct scenario $s_i$ therefore
contributes at least one case whose correct output is determined by
$\phi_i$, and the precedence ordering is witnessed by all $n$ cases jointly.
\end{proof}

Precedence defects concentrate near guard boundaries---the thin set where
one scenario yields to another.
Random tests rarely land there; when they miss, unit-test repair looks
healthy while ordering bugs remain.

\begin{definition}[Boundary Region]
\label{def:boundary-region}
For margin $\varepsilon > 0$, the \emph{boundary region} is
$\mathcal{B}_\varepsilon(t) = \{\mathbf{x} \in \mathcal{X}_t :
  \exists i,\; \partial g_i(\mathbf{x}) < \varepsilon\}$.
\end{definition}

Fault-exposure probability is how often a faulty candidate fails on a
randomly drawn witness; when it is small, denser unit tests still look fine.

\begin{definition}[Per-witness Fault-Exposure Probability]
\label{def:fault-exposure}
For faulty candidate $c$, $\varepsilon_c =
\Pr_{\mathbf{x} \sim \mathcal{D}(t)}[\,f_c(\mathbf{x}) \neq g_t(\mathbf{x})\,]$;
$\varepsilon = \min_{c \in \mathcal{C}\setminus\{g_t\}} \varepsilon_c > 0$.
\end{definition}

The corollary makes the random-vs-SMT contrast quantitative.

\begin{corollary}[Random vs.\ SMT Fault-Exposure Hit Rate]
\label{cor:random-vs-smt}
Let $\varepsilon = \mu(\mathcal{B}_\varepsilon(t))/\mu(\mathcal{X}_t)$ denote the
measure of the boundary region relative to the input domain.
For a faulty candidate $c$ with per-witness fault-exposure probability
$\varepsilon_c \geq \varepsilon$ (Definition~\ref{def:fault-exposure}),
a single uniform random sample detects the fault with probability at most
$\varepsilon_c$, because
$\Pr[\mathbf{x} \in \mathcal{B}_\varepsilon(t)] =
\mu(\mathcal{B}_\varepsilon(t))/\mu(\mathcal{X}_t)$ under uniform sampling---
typically small when guard overlap is dense.
The Constraint-guided Witness Generator produces one witness per strict
scenario region via $\Phi_i$ (Equation~\eqref{eq:priority}); under sound
encoding, each such witness activates scenario~$i$ with probability~$1$
(Theorem~\ref{thm:witness-soundness}), lying on (or within numeric precision of)
the boundary of~$g_i$.
Random sampling therefore reaches precedence-critical regions with probability
$\mathcal{O}(\varepsilon)$, whereas SMT achieves targeted boundary coverage
deterministically.
\end{corollary}

\begin{proof}
Assume faults of interest manifest only on inputs in
$\mathcal{B}_\varepsilon(t)$.
Under uniform sampling,
$\Pr[f_c(\mathbf{x})\neq g_t(\mathbf{x})]
  \le \Pr[\mathbf{x}\in\mathcal{B}_\varepsilon(t)]=\varepsilon$.
For SMT, satisfiability of $\Phi_i$ yields
$g_i(\mathbf{x}_i^*)\land\bigwedge_{j<i}\neg g_j(\mathbf{x}_i^*)$, so
$\mathbf{x}_i^*$ lies in the exclusive first-match region of~$s_i$ and
activates that region with probability~$1$ under a sound encoding
(Theorem~\ref{thm:witness-soundness}).
Hence random sampling hits precedence-critical regions with probability
$\mathcal{O}(\varepsilon)$, while SMT hits them deterministically.
\end{proof}

\begin{figure}[t]
  \centering
  \tikzinput{diagrams/tikz/fsf_semantics}
  \caption{\textbf{FSF evaluation semantics for \code{classify\_signal}.}
    Left: the three-scenario FSF definition.
    Right: input-space partitioning with priority $S_1 \succ S_2 \succ S_3$;
    the green marked point $(-3,-10)$ is an $S_1$ witness that also lies in
    the geometric region of $S_2$'s guard.
    Takeaway: first-match, not geometric inclusion alone, decides the oracle.}
  \label{fig:fsf-semantics}
\end{figure}

The point of Figure~\ref{fig:fsf-semantics} is the green witness $(-3,-10)$:
geometrically it also satisfies $S_2$'s guard, yet first-match assigns it to
$S_1$.
SMT case generation places one such witness in each strict scenario region.

Pass/fail alone is not enough for repair.
The pipeline also needs a counterexample that names the violated scenario,
so the model can localise the bug inside the ordered specification.

\begin{definition}[Strict Scenario Violation]
\label{def:violation}
A candidate $c$ \emph{violates scenario $s_i$} on the concrete case
$(x_i, y_i) \in \mathcal{D}(t)$ if
\[
  f_c(x_i) \;\neq\; y_i \;=\; g_t(x_i),
\]
where $f_c : \mathcal{X}_t \to \mathcal{Y}_t$ denotes the observable
behaviour of $c$.  The associated \emph{counterexample} is the triple
\[
  (x_i,\;g_t(x_i),\;f_c(x_i)),
\]
comprising the failing input, the expected output, and the actual output
produced by $c$.
\end{definition}

\paragraph{Example.}
For Listing~\ref{lst:llm-incorrect} on $(-3,-10)$, the counterexample is
$\bigl((-3,-10),\;\texttt{Error},\;\texttt{Critical}\bigr)$ and the violated
scenario is $s_1$.

Counterexamples are the primary feedback signal returned to the LLM.
They are structured: the pipeline identifies not only that the candidate
produced a wrong answer, but \emph{which scenario} was active, allowing the
language model to localise the bug within the FSF specification.

% ---------------------------------------------------------------------------
\section{The Bounded Acceptance Problem}
\label{sec:form:acceptance}
% ---------------------------------------------------------------------------

The oracle and the case set answer whether a candidate is right on the
witnesses that matter.
Release is stronger: clear those witnesses \emph{and} a structural screen,
within a fixed LLM budget.

\begin{definition}[Candidate Implementation]
\label{def:candidate}
A \emph{candidate implementation} for task $t$ is any computable function
$c : \mathcal{X}_t \to \mathcal{Y}_t$.
Write $f_c$ for its observable input-output behaviour.
\end{definition}

The natural quality score is the fraction of witnesses on which $c$ matches
the oracle---continuous, so repair progress remains visible, and grounded in
the specification rather than an arbitrary unit-test suite.

\begin{definition}[Strict Formal Conformance]
\label{def:conformance}
\[
  \mathrm{Conf}(c,\,t)
  \;=\;
  \frac{1}{|\mathcal{D}(t)|}
  \sum_{(x,\,y)\,\in\,\mathcal{D}(t)}
  \mathbf{1}\bigl[f_c(x) = y\bigr]
  \;\in\; [0,1].
\]
\end{definition}

\paragraph{Example.}
If $\mathcal{D}(t)$ has three witnesses and the buggy
\code{classify\_signal} fails only on $(-3,-10)$, then
$\mathrm{Conf}(c,t)=2/3$.

\emph{Strict success} is the thresholded form: $\mathrm{Conf}(c,t)\ge\tau$
with $\tau=1$ (Definition~\ref{def:strict-success}).
A lower $\tau$ would mask ordering violations that pass most cases where
overlapping guards prescribe identical outputs~\citep{zhu1997specification}.

\begin{definition}[Strict Success]
\label{def:strict-success}
Candidate $c$ \emph{strictly succeeds} on task $t$ iff
$\mathrm{Conf}(c,\,t) \ge \tau$ with $\tau = 1$.
\end{definition}

Conformance alone is still insufficient: structural anti-patterns (hard-coded
returns, unconditional success) can clear every witness.
The high-severity screen $\mathcal{P}^H$ (nine patterns in
Table~\ref{tab:patterns}; Appendix~\ref{app:patterns}) blocks those cases;
medium/low patterns inform repair without gating release.

\begin{definition}[High-Severity Pattern Violation]
\label{def:pattern-violation}
Candidate $c$ is \emph{pattern-safe} iff
$\forall p \in \mathcal{P}^H : p(c,\,t) = 0$, where each checker
$p : (c, t) \mapsto \{0, 1\}$ returns $1$ on a hit.
\end{definition}

Release is the conjunction of the two checks---neither alone is enough.

\begin{definition}[Hybrid Acceptance Predicate]
\label{def:accept}
\[
  \mathrm{Accept}(c,\,t)
  \;=\;
  \mathbf{1}\!\left[
    \mathrm{Conf}(c,\,t) \geq 1
    \;\land\;
    \forall p \in \mathcal{P}^H : p(c,\,t) = 0
  \right].
\]
\end{definition}

\paragraph{Example.}
A repaired \code{classify\_signal} with $\mathrm{Conf}=1$ still fails
acceptance if RF07 (\emph{unconditional success}) fires on a hard-coded
\texttt{return "Normal"}.
Conversely, clearing the pattern screen while failing one witness also fails
$\mathrm{Accept}$.
That conjunction is the source of the strict-success paradox in
Chapter~\ref{ch:results}: mean conformance can be high while acceptance stays
selective.

Those pieces assemble into the computational problem: find an acceptable
candidate within budget~$K$, or return the best-effort candidate when the
budget is exhausted.
The feedback object on attempt~$k{+}1$ is a structured JSON record
(counterexamples, pattern hits, attempt index); Chapter~\ref{ch:method}
details its fields.
The search is CEGIS-shaped~\citep{solarlezama2008sketching,lynn2024verifierloop}
with a stochastic synthesiser---hence the explicit fallback.

\begin{definition}[Bounded Acceptance Problem]
\label{def:bap}
Given task $t$, budget $K \in \mathbb{N}^+$, and synthesis oracle
$\mathrm{LLM}(\cdot)$, find
$c^* = \mathrm{LLM}(t,\mathrm{feedback}_1,\ldots,\mathrm{feedback}_{k-1})$
with at most $K$ calls such that $\mathrm{Accept}(c^*, t) = 1$;
otherwise return
$c^* = \arg\max_{c \in C_K} \mathrm{Conf}(c,\,t)$
where $C_K$ is the set of candidates generated in the $K$ attempts.
\end{definition}

% ---------------------------------------------------------------------------
\section{Prevention and Quality Metrics}
\label{sec:form:metrics}
% ---------------------------------------------------------------------------

Generation success is the wrong primary metric for a prevention system.
What matters is whether faulty candidates are kept out of release, and whether
correct ones are not rejected by accident.
Both rates are measured over a mutant population
\citep{demillo1978mutation,ammann2016testing}.

\begin{definition}[Mutant Population]
\label{def:mutants}
A \emph{mutant population} $\mathcal{M}$ is a finite set of pairs
$(c_m, t_m)$ where at least one of the two has been artificially perturbed.
\emph{Specification-confusion} mutants $\mathcal{M}^S$ mutate the task
(guards or outputs) while keeping a correct implementation of the original;
\emph{implementation-screening} mutants $\mathcal{M}^I$ mutate the candidate
(e.g.\ guard-order swap) while keeping the task.
Write $\mathrm{Faulty}(m)=1$ when the pair is genuinely faulty.
\end{definition}

PDR is recall-like: of all mutants, how many does the pipeline refuse?
FAR is precision-like: of genuinely faulty mutants, how many still slip
through $\mathrm{Accept}$?
In a sprint review, high PDR and low FAR mean the gate is selective without
rubber-stamping~\citep{okun2007mutation}.

\begin{definition}[Prevention Detection Rate]
\label{def:pdr}
\[
  \mathrm{PDR}
  \;=\;
  \frac{\bigl|\bigl\{m \in \mathcal{M} : \neg\,\mathrm{Accept}(c_m,\,t_m)\bigr\}\bigr|}
       {|\mathcal{M}|}.
\]
\end{definition}

FAR is the complementary slip-through rate.

\begin{definition}[False Accept Rate]
\label{def:far}
\[
  \mathrm{FAR}
  \;=\;
  \frac{\bigl|\bigl\{m \in \mathcal{M} :
      \mathrm{Accept}(c_m,\,t_m) = 1
      \;\land\; \mathrm{Faulty}(m) = 1
    \bigr\}\bigr|}
       {|\mathcal{M}|}.
\]
\end{definition}

Mutation kill rate $\mu_k$ isolates the formal check alone
(Definition~\ref{def:mkr}): the fraction of mutants with
$\mathrm{Conf}<1$.
The gap $\mathrm{PDR}-\mu_k$ is the additive value of pattern screening---mutants
that clear every witness yet still fail $\mathrm{Accept}$.

\begin{definition}[Mutation Kill Rate]
\label{def:mkr}
\[
  \mu_k
  \;=\;
  \frac{\bigl|\bigl\{m \in \mathcal{M} : \mathrm{Conf}(c_m,\,t_m) < 1\bigr\}\bigr|}
       {|\mathcal{M}|}.
\]
\end{definition}

% ---------------------------------------------------------------------------
\section{Benchmark Formalization}
\label{sec:form:benchmark}
% ---------------------------------------------------------------------------

The evaluation uses a synthetic benchmark of hard tasks built to stress FSF
precedence sensitivity \citep{li2023iet}.
The construction below makes that stress precise and proves that correct
solutions must respect scenario ordering.

\begin{definition}[Precedence-Sensitive Task]
\label{def:prec-sensitive}
A task $t$ is \emph{precedence-sensitive} if there exist two scenarios
$s_i, s_j$ with $i < j$ such that their guard regions overlap:
\[
  \bigl\{x \in \mathcal{X}_t : g_i(x) = 1\bigr\}
  \;\cap\;
  \bigl\{x \in \mathcal{X}_t : g_j(x) = 1\bigr\}
  \;\neq\; \emptyset.
\]
\end{definition}

In a precedence-sensitive task, the output on inputs in the overlapping region
depends entirely on which guard is evaluated first.  An implementation that
checks $g_j$ before $g_i$ will, on those inputs, return $\phi_j(x)$ rather
than the correct $\phi_i(x)$.  This is precisely the class of error that
naive LLM synthesis frequently introduces: the model may faithfully reproduce
all guard conditions and output mappings from the specification text, while
inadvertently evaluating them in a different order.

\begin{proposition}
\label{prop:order-required}
For any precedence-sensitive task $t^H$, a strictly correct implementation
must evaluate the scenario guards in the exact order prescribed by
$\mathcal{S}_{t^H}$.
\end{proposition}

\begin{proof}
Let $x^*$ be any input in the overlapping region of $s_i$ and $s_j$
($i < j$).  By Definition~\ref{def:oracle}, the oracle maps
$g_t(x^*) = \phi_i(x^*)$.  Suppose an implementation $c$ evaluates $g_j$
before $g_i$.  Since $g_j(x^*) = 1$, the implementation returns
$\phi_j(x^*)$.  Because the task is precedence-sensitive (the two scenarios
cover an overlapping region of the input space), the specification is
non-trivially ordered, and it need not be the case that $\phi_i(x^*) =
\phi_j(x^*)$ — indeed, if they were always equal on the overlap, the overlap
would be semantically irrelevant and the task would not test ordering at all.
Hence $f_c(x^*) \neq g_t(x^*)$, and $c$ does not strictly succeed.
\end{proof}

Hard synthetic tasks are constructed as follows.
Each task $t^H$ has the following structure:

\begin{itemize}
  \item \textbf{Input variables}: five non-negative integers
        $a, b, c, d, e \in \mathbb{N}$.
  \item \textbf{Output variables}: three integer quantities —
        $\mathit{risk} \in \{0,1,2\}$,
        $\mathit{action} \in \{0,1,2\}$, and
        $\mathit{score} \in [0, 100]$ — giving $\mathcal{Y}_t \subseteq
        \{0,1,2\}^2 \times [0,100]$.
  \item \textbf{Scenario count}: up to $n = 9$ ordered scenarios.
  \item \textbf{Threshold parameters}: seven threshold values
        $\theta_1, \ldots, \theta_7$ sampled uniformly at random from
        $[1, 100] \cap \mathbb{Z}$, used to construct the guard predicates
        (e.g., $a > \theta_1$, $b \leq \theta_2 \land c > \theta_3$).
  \item \textbf{Feasibility guarantee}: all scenarios are satisfiable, verified
        by a Z3 feasibility check during task generation.  Any candidate
        $\theta$ vector that renders a scenario unsatisfiable is discarded
        and resampled.
  \item \textbf{Consistency validation}: the generated task passes a
        scenario-consistency validator that confirms each generated case
        belongs to exactly one scenario and that the ``others'' guard
        correctly covers all remaining inputs.
\end{itemize}

By Proposition~\ref{prop:order-required}, every hard task in the benchmark
requires an implementation to evaluate guards in the specified order.  An
LLM that does not attend to the ordering of scenarios in the prompt will
therefore produce incorrect implementations, which are caught by the
SMT-generated test suite.

% ---------------------------------------------------------------------------
\section{Computational Considerations}
\label{sec:form:complexity}
% ---------------------------------------------------------------------------

Formal checking must not erase the prevention benefit by dominating runtime.
The cost accounting below confirms that Z3 and pattern screening remain cheap
relative to LLM inference \citep{demoura2008z3, barrett2018smt}.

\paragraph{Case Generation.}
For each of the $n$ scenarios, the Z3 solver is invoked with a linear
arithmetic formula over $d$ integer variables.  Linear integer arithmetic
(LIA) is decidable, and Z3's LIA solver runs in time that is polynomial in
the formula size for the types of constraints that arise from threshold-based
FSF guards.  In practice, with $n \leq 9$ and $d = 5$, each call to Z3
completes in under $0.1$ seconds, giving a total case generation time of
$O(n \cdot d^2)$ solver steps in the worst case, or approximately
$0.5$--$1$ seconds per task.

\paragraph{Conformance Checking.}
Given the case set $\mathcal{D}(t)$ of size $n$, checking conformance
requires executing the candidate $c$ on each input $x_i$ and comparing the
output to $g_t(x_i)$.  This is a simple lookup-and-compare operation with
cost $O(n)$ function evaluations and $O(n)$ comparisons.  For $n \leq 9$,
this step is effectively instantaneous.

\paragraph{Pattern Matching.}
The pattern screening layer applies $|\mathcal{P}^H| = 14$ pattern checkers
to the source code of candidate $c$.  Each pattern is encoded as a regular
expression or a simple abstract-syntax-tree traversal.  The cost is $O(L
\cdot P)$ where $L$ is the length of the candidate source code and $P =
|\mathcal{P}^H|$.  For typical LLM-generated functions of 30--80 lines,
this step completes in well under $0.01$ seconds.

\paragraph{Total Pipeline Cost.}
The dominant cost at every iteration is the LLM API call, which takes
approximately $10$ seconds per attempt in the reported experiments.
Z3 solving and conformance checking add less than $1$ second combined;
pattern matching is negligible.
With budget $K = 3$, expected runtime is about $30$ seconds per task, of
which more than $97\%$ is LLM inference.
Formal checking is therefore not the bottleneck: larger case sets or richer
pattern libraries would not meaningfully change throughput.

\paragraph{Summary.}
The cost structure can be written as
\[
  T_{\mathrm{pipeline}}
  \;\approx\;
  K \cdot T_{\mathrm{LLM}}
  \;+\; O(n \cdot d^2) \cdot T_{\mathrm{Z3}}
  \;+\; O(n) \cdot T_{\mathrm{exec}}
  \;+\; O(L \cdot P) \cdot T_{\mathrm{pattern}},
\]
where $T_{\mathrm{LLM}} \gg T_{\mathrm{Z3}} \gg T_{\mathrm{exec}} \approx
T_{\mathrm{pattern}}$, confirming that formal checking adds less overhead
than a single LLM call.

% ---------------------------------------------------------------------------
\section*{Chapter Summary}
% ---------------------------------------------------------------------------

What must be true is now fixed: ordered tasks, SpecIR as the shared hinge,
first-match oracles, SMT witnesses with full scenario coverage
(Lemma~\ref{lem:coverage}), and $\mathrm{Accept}$ under budget~$K$
(Definition~\ref{def:bap}).
PDR/FAR measure prevention; every hard task requires ordering
(Proposition~\ref{prop:order-required}); formal checking stays cheap relative
to LLM inference.
Chapter~\ref{ch:method} turns these objects into a runnable loop.
